\documentclass[11pt]{article}
\usepackage[a4paper,margin=1in]{geometry}
\usepackage{amsmath,amssymb,amsthm}
\title{Sharp Summary: Maas (2011) -- Entropy Gradient Flows for Finite Markov Chains}
\author{}
\date{}
\begin{document}
\maketitle

\section*{Problem and scope}
Given an irreducible, reversible Markov kernel $K$ on a finite set $\mathcal X$ (with invariant measure $\pi$), the paper asks:
\emph{can the Markov heat flow be written as a gradient flow of entropy, as in JKO/Wasserstein theory in $\mathbb R^n$?}

This is not possible with classical $W_2$ on discrete spaces. The paper constructs a new transport metric $\mathcal W$ that restores this structure.

\section*{Core construction}
For densities
\[
\mathcal P(\mathcal X)=\Big\{\rho\ge 0:\sum_x \pi(x)\rho(x)=1\Big\},
\qquad
\mathcal H(\rho)=\sum_x \pi(x)\rho(x)\log\rho(x),
\]
choose a mean function $\theta$ (logarithmic mean is the key case), and define
\[
\rho(x,y)=\theta(\rho(x),\rho(y)).
\]
Then define $\mathcal W$ by a discrete Benamou--Brenier formula:
\[
\mathcal W(\rho_0,\rho_1)^2
=\inf_{\rho_t,\psi_t}
\frac12\int_0^1\sum_{x,y}\big(\psi_t(x)-\psi_t(y)\big)^2 K(x,y)\rho_t(x,y)\pi(x)\,dt,
\]
subject to the discrete continuity equation
\[
\frac{d}{dt}\rho_t(x)
+\sum_y\big(\psi_t(y)-\psi_t(x)\big)K(x,y)\rho_t(x,y)=0.
\]

So transport cost is edge-weighted by both graph geometry ($K$) and current mass ($\rho$).

\section*{Main theorem (precise statement level)}
Let $H(t)=e^{t(K-I)}$ be the heat semigroup. With the logarithmic mean choice for $\theta$, the paper proves:
\[
\rho_t = H(t)\rho
\quad\Longrightarrow\quad
D_t\rho_t = -\nabla_{\mathcal W}\mathcal H(\rho_t),
\]
i.e., heat flow is the gradient flow of relative entropy in the Riemannian structure induced by $\mathcal W$.

It also characterizes finiteness/topology properties of $\mathcal W$ and geodesic structure on components of $\mathcal P(\mathcal X)$.

\section*{Insight for discrete diffusion}
This paper does \textbf{not} give a DDPM-style learning objective or reverse-kernel estimator. Its contribution is foundational geometry:

\begin{itemize}
\item Forward Markov noising can be viewed as entropy-dissipating transport.
\item The right metric in discrete space is not classical Wasserstein, but a Markov-kernel-adapted metric.
\item This suggests principled regularizers for learned reverse kernels $Q^t$ (entropy production, geometry-aware smoothness), rather than ad-hoc constraints.
\end{itemize}

\section*{Relevance to this repo}
\textbf{Indirect but high-value theory.}
It is not a discrete diffusion algorithm paper, but it provides the most precise mathematical lens for the ``transport'' interpretation mentioned in the tutorial and internship text.

\end{document}
